\section{Limitations and conclusion}
\label{sec:conclusion}

%% ===========================================================================
%% SKELETON — write this yourself, as with the introduction.
%%
%% Rubric: minimum 300 words, and limitations are graded explicitly. The top
%% band asks for major limitations discussed in depth, so this section rewards
%% candour rather than defensiveness. Every item below is a real limitation of
%% this work, not a hypothetical one.
%% ===========================================================================

%% ---------------------------------------------------------------------------
%% LIMITATIONS  (~180 words)
%%
%%  - Corpus size. Small relative to what a fine-tuned transformer can absorb.
%%    Quantify the consequence rather than merely admitting it: the 95% confidence
%%    interval on F1 is +/-0.067 at 50 test sentences, +/-0.050 at 100, +/-0.036
%%    at 250. State the test split size and the resulting interval, and say plainly
%%    what magnitude of difference the study can and cannot detect.
%%
%%  - Domain. Wikipedia summaries only. Encyclopaedic register, entity-dense,
%%    unrepresentative of news, speech, or social media.
%%
%%  - Single adjudicator. Disagreements were resolved by one person, so
%%    systematic bias in that judgement is not detectable by the agreement
%%    statistics, which measure annotators against each other and not against
%%    truth.
%%
%%  - Annotator anchoring. Even in the from-scratch condition, annotators worked
%%    from the same written guidelines, which were themselves informed by model
%%    output. The conditions differ in whether suggestions were visible, not in
%%    whether the model influenced the standard.
%%
%%  - Annotator protocol compliance. One of ten annotators (id 11, batch 12) never
%%    marked a head: 0 of 73 multi-token entities, where peers narrow 10-49% of the
%%    time. Under the lowest peer rate the chance of that being coincidence is 3e-4.
%%    Their 100 sentences are in the gold set with heads defaulting to full spans, so
%%    they contribute 11.5% of the multi-token pool with a systematic bias toward
%%    full-span behaviour. Say what you did about it -- excluded them from head-view
%%    results, had the batch re-annotated, or reported the head view with the caveat.
%%    See results/annotator_qc.json. This is worth reporting rather than quietly
%%    fixing: it is evidence that a two-boundary protocol costs compliance, which is a
%%    finding about the method, not just a defect in one annotator's work.
%%
%%  - Automatic adjudication. The 20 shared sentences were resolved by token-level
%%    majority vote, not by a human. 12 of the 20 were non-unanimous
%%    (results/adjudication.md). Majority vote is reproducible but it is not
%%    adjudication, and the two should not be conflated in the write-up.
%%
%%  - Reproducibility. Training was not bit-reproducible on the hardware used:
%%    identical seeds gave F1 0.9187 and 0.9279 (Section~\ref{sec:method:stats}).
%%    Numbers are reproducible in distribution, not exactly.
%%
%%  - Statistical power for active learning. The detection threshold of roughly
%%    0.04 (Section~\ref{sec:results:al}) means a null result here does not show
%%    that acquisition strategy is irrelevant, only that any effect is smaller
%%    than this design can resolve.
%% ---------------------------------------------------------------------------

%% ---------------------------------------------------------------------------
%% CONCLUSION  (~120 words)
%%
%%  Restate what was built and what was measured, and let the results carry it.
%%  Do not introduce claims the results section does not support.
%%
%%  Candidate future work, only if genuinely warranted:
%%    - extend beyond Wikipedia to a second domain
%%    - a second adjudicator on a subset, to bound adjudication bias
%%    - a wider budget range for the active learning comparison, since the
%%      present design cannot resolve small effects
%% ---------------------------------------------------------------------------
